\bigskip\Needspace{10\baselineskip}
\pdfbookmark[1]{Fortgesetzte Nachricht. Ceres. Februar 1802}{ceres180202}
\begin{center}
{\small FORTGESETZTE NACHRICHT ÜBER DEN NEUEN HAUPTPLANETEN.\par}
[CERES.]\par
\medskip\hrule\medskip
Monatliche Correspondenz. Februar 1802.\par
\medskip\hrule
\end{center}\medskip
—\quad —\quad —\quad „Die Übereinstimmung der Deklinationen in den Beobachtungen \textsc{Piazzi}s sowohl untereinander als auch mit der Rechnung ist tatsächlich bewundernswert. Gerade das bestärkt meine Hoffnung, dass meine Elemente für die Auffindung genau genug sein werden. Denn gerade die Breitenbeobachtungen, die am stärksten von den Deklinationen abhängen, haben den größten Einfluss auf die Übereinstimmung der Elemente. Eine Änderung um $20''$ in der Breite würde die Ellipse völlig verändern, wenn man sie bei den äußeren Beobachtungen positiv und bei den mittleren negativ anbrächte oder umgekehrt. Damit hängt eine weitere Bemerkung zusammen, die aber erst dann ihren vollen Wert erhält, wenn die Wiederauffindung gelungen ist und die aus den Beobachtungen abgeleiteten Elemente bestätigt hat. In diesem Fall werden wir die Wiederauffindung hauptsächlich der starken Neigung der Bahn verdanken. Fiele die Bahn mit der Ekliptik zusammen, so müsste man wohl auf die Wiederauffindung verzichten. Zumindest wäre die Ungewissheit in der Länge nach einem Jahr außerordentlich groß und betrüge vielleicht ein paar Zeichen, obwohl diese Zone dann fast gar keine Breite hätte. Der Grund ist leicht einzusehen. Dann wären zur Bestimmung der Ellipse zwingend vier Längen erforderlich. Zwei äußere und eine mittlere Beobachtung ließen sich folglich durch unendlich viele Ellipsen und vielleicht auch durch Parabeln und Hyperbeln vollkommen darstellen. Bei der Kürze der Beobachtungsreihe ist leicht abzusehen,
\SourceLeaf{210}%
dass Elemente, die zwei äußere und eine mittlere Beobachtung darstellen, auch von der ganzen Reihe nur wenig abweichen können. Der Kegelschnitt wäre folglich fast ganz unbestimmt...“

Wie richtig Dr.\ \textsc{Gauss} dies beurteilt hat, liegt nun bei der Wiederauffindung der \textit{Ceres} zutage. Die enge Übereinstimmung seiner Elemente mit den jetzigen Beobachtungen bestätigt es. Ja, ohne die scharfsinnigen Bemühungen und Berechnungen von Dr.\ \textsc{Gauss} hätten wir die \textit{Ceres} vielleicht nicht wiedergefunden. Der größere und schönere Teil des Verdienstes gebührt daher ihm. Dieser Meinung ist auch der verehrte Dr.\ \textsc{Olbers}. Er schreibt an den Herausgeber: „Mit Vergnügen werden Sie bemerkt haben, wie genau die Ellipse von Dr.\ \textsc{Gauss} mit den Beobachtungen der \textit{Ceres} übereinstimmt. Teilen Sie dies doch diesem verehrten Gelehrten mit und übermitteln Sie ihm dabei meine ganz besondere Hochachtung. Ohne seine mühsamen Untersuchungen über die elliptischen Elemente dieses Planeten hätten wir ihn vielleicht überhaupt nicht wiedergefunden. Ich jedenfalls hätte ihn nicht so weit östlich gesucht.“
\begin{flushright}[\textsc{von Zach}.]\end{flushright}

